% Generated draft: requires mathematical and source review.
\documentclass[11pt]{article}
\usepackage[margin=25mm]{geometry}
\usepackage{fontspec}
\setmainfont{Latin Modern Roman}
\usepackage{amsmath,amssymb,mathtools,unicode-math}
\setmathfont{Latin Modern Math}
\usepackage{enumitem,needspace,xurl,hyperref}
\setlength{\parindent}{0pt}
\setlength{\parskip}{5pt}
\setlength{\emergencystretch}{3em}
\newcommand{\N}{\mathbb{N}}
\newcommand{\Z}{\mathbb{Z}}
\newcommand{\Q}{\mathbb{Q}}
\newcommand{\R}{\mathbb{R}}
\newcommand{\C}{\mathbb{C}}
\newcommand{\F}{\mathbb{F}}
\newcommand{\A}{\mathbb{A}}
\newcommand{\PP}{\mathbb{P}}
\newcommand{\E}{\mathbb{E}}
\newcommand{\eps}{\varepsilon}
\newcommand{\dd}{\,\mathrm{d}}
\begin{document}
\section*{UnsolvedMath ID: 30005140}
\textbf{Source problem code:} catalogue entry
\subsection{Definitions and mathematical statement}
Let \(p\in(1/2,1]\). In Bernoulli bond percolation, each edge of the underlying lattice is independently declared open with probability \(p\). Write \(\omega\) for the percolation configuration and \(\mathcal C_\infty(\omega)\) for its infinite open cluster. The supplied statement concerns two-dimensional percolation; the threshold \(1/2\) suggests the nearest-neighbor square lattice \(\Z^2\), but the lattice is not explicitly specified in the record.

For a specified finite domain and boundary convention, a discrete Gaussian free field is a centered Gaussian field whose covariance is the corresponding Green function, with the Laplacian and Green-function normalization fixed. The record does not specify the finite-volume construction, boundary conditions, or normalization of the field on the infinite cluster.

Let \(B_N\) denote the intended lattice box of side length \(N\), and let \(h_N\) denote the intended finite-volume field associated with \(\mathcal C_\infty(\omega)\) and \(B_N\). Subject to those missing specifications, write
\[
M_N=\max_{x\in B_N\cap\mathcal C_\infty(\omega)}h_N(x).
\]
The box placement and the convention when the indexing set is empty must also be fixed. Write \(a_N\) for the centering prescribed by the result being extended; its formula and any dependence on \(p\) or \(\omega\) are not supplied.

\textbf{Intended question.} For every fixed \(p\in(1/2,1]\), does the centered maximum satisfy
\[
M_N-a_N\xrightarrow[N\to\infty]{\mathrm d}L_p
\]
for a limiting probability law, under the field construction, centering, and probability-law convention of the cited result? The notation \(L_p\) is only a placeholder: the record does not determine whether the intended convergence is conditional on the environment or averaged over it, or whether a particular limiting law is required. The centering must be recovered from the cited result and is not a freely chosen sequence.

The supplied dated assessment states that this convergence has been proved for \(p\) sufficiently close to \(1\), and asks for its extension to every \(p>1/2\). That assertion is supplied curation, not an independently verified mathematical or literature conclusion. Because the excerpt refers to an unspecified preceding result, the record does not determine a unique fully specified convergence statement.
\subsection{Short English statement}
The question asks whether the centered maximum of the discrete Gaussian free field on a two-dimensional infinite percolation cluster converges in distribution for every fixed \(p>1/2\), extending the range near \(1\) described in the supplied assessment. The field construction, boundary conditions, centering, and conditional or averaged interpretation of convergence must be recovered before this becomes a fully precise statement.
\subsection{Sources}
\begin{itemize}
\item Local UnsolvedMath registry, permanent ID \(30005140\), title: Gaussian-Free-Field Maxima on Percolation Clusters; exact source problem code: \texttt{catalogue entry}. The record belongs to the Oberwolfach Reports Open Problems collection and supplies proposed year \(2022\); no individual proposer is identified.
\item Supplied source citation: \emph{Universality: Random Matrices, Random Geometry and SPDEs} (2023), Oberwolfach Reports, pp.\ 1501--1550, DOI \texttt{10.4171/owr/2022/27}. \url{https://doi.org/10.4171/owr/2022/27}
\item Supplied literature reference: Schweiger--Zeitouni, DGFF maximum. \url{https://doi.org/10.1002/cpa.22181} The supplied record does not provide the full bibliographic title or the relevant theorem statement.
\item Supplied literature reference: Oberwolfach report. \url{https://doi.org/10.4171/owr/2022/27}
\item Literature assessment dated \texttt{2026-08-22}: the record labels the problem open and states that the centered maximum limit is proved for \(p\) sufficiently close to \(1\), with the extension to all \(p>1/2\) remaining open. This is supplied dated curation derived from the OWR material, not independently verified here; its status and cited claims require independent checking.
\item The registry labels the statement \texttt{corrected\_verified} and reports that the official Oberwolfach report was checked directly. This is a supplied verification claim, not a source visit or verification performed for this draft.
\end{itemize}
\end{document}
